\section{A checkable matching bound for restricted graph squares}
\label{sec:matching-bound}
Square-clique certificates can be replaced by a chromatic lower bound
when a minimum square coloring has a checkable matching representation.
The following argument applies classical matching theory; it is not a
claim that the Tutte--Berge theorem or this elementary coloring identity
is new.

\begin{proposition}
\label{prop:matching-bound}
Let $G$ be a nonempty simple undirected graph with $n$ vertices and
$h\geq k\geq1$. Write $H=\overline{G^2}$. If $H$ is triangle-free, then
\[
 \chi(G^2)=n-\nu(H),\qquad
 \lambda_{h,k}(G)\geq k\bigl(n-\nu(H)-1\bigr),
\]
where $\nu(H)$ is the maximum matching size. A matching $M$ and a set
$S\subseteq V(H)$ satisfying
\begin{equation}
 2|M|=n+|S|-o(H-S)
 \label{eq:matching-certificate}
\end{equation}
certify that $\nu(H)=|M|$, where $o(H-S)$ counts odd-order components.
If $h=k$, matched pairs and unmatched singletons yield an optimal
labeling with spacing $k$. If $k<h\leq2k$, the same conclusion holds
whenever these classes admit a compatible order as in
Proposition~\ref{prop:color-order}.
\end{proposition}
\begin{proof}
Normalize any feasible labeling so its smallest label is zero, and
write $s$ for its span. The bins $\lfloor f(v)/k\rfloor$ form a proper
coloring of $G^2$: two labels in the same bin differ by less than $k$,
which is forbidden at distance one or two. Hence
$\chi(G^2)\leq\lfloor s/k\rfloor+1$, giving
$s\geq k(\chi(G^2)-1)$.

Since $H$ is triangle-free, every independent set of $G^2$ has at most
two vertices. In any coloring, its two-vertex classes form a matching
in $H$. Conversely, every matching in $H$, completed by unmatched
singletons, defines a coloring of $G^2$. A coloring with $t$ pairs uses
$n-t$ colors, proving the identity for $\chi(G^2)$.

For any matching $N$ in $H$, at most $|S|$ of its edges touch $S$.
In each component of $H-S$, the remaining edges cover at most the
largest even number of vertices in that component. Therefore
$2|N|\leq n+|S|-o(H-S)$. Equality for a valid $M$ proves it maximum.
Finally, spacing the classes by $k$ is valid for $h=k$; for $h>k$,
the compatible-order condition proves validity. In either case the
span attains the established lower bound.
\end{proof}

The existence of a set $S$ satisfying~\eqref{eq:matching-certificate}
for a maximum matching follows from the Tutte--Berge formula
\cite{goemans2017matching}. The generator first checks $S=\varnothing$.
If it is insufficient, it forms
$D=\{v:\nu(H-v)=\nu(H)\}$ and uses $S=N_H(D)\setminus D$, the
Gallai--Edmonds barrier. It makes at most $n+1$ matching calls. The
standalone checker reconstructs $H$ from original edges, checks its
triangle-free condition, validates the matching and counts components
after deleting $S$. Verification uses neither NetworkX nor a solver's
claim that a matching is maximum. Optional labelings are checked
directly; optimality is reported only if the two bounds meet.

The chromatic bound can exceed the clique bound: $C_7^2$ has clique
number three and chromatic number four, with a seven-cycle complement.
Thus the matching bound is $3k$ whereas the clique bound is $2k$.
This example does not imply improvement over every other graph bound.
In the actual original cohort, the condition holds for 15/117 instances;
all 15 already have equally strong certificates. No bound is improved
and no additional interval closes. The independent-support total stays
90, with 11 other closed cases relying on saved solver lower bounds.
On the new cohort the condition holds for 39/207 instances; again it
never strengthens the common degree/clique bound. These null results
are retained. The matching certificate expands the types of evidence
the checker accepts, but provides no empirical coverage gain in these
two cohorts. The earlier certificate counts and timed measurements are
unchanged.
